\subsubsection*{The smooth-curve step in cyclotomic reconstruction}

Fix an algebraic subfield $k\subset\overline{\mathbb Q}$ for which
$k\subset F\mathbb Q^{\mathrm{ab}}$ for some number field $F$.
Every finite extension of $k$ has the same property, after enlarging
$F$. Write $K=\overline{\mathbb Q}$. We use the rational-source
integrality lemma proved above, with all multiplicative completions
taken as $\varprojlim_n A/A^n$.

\emph{Arithmetic input on abelian varieties.}
For each finite extension $L/k$ and each abelian variety $A/L$,
\[
 \bigcap_{n\geq1} nA(L)=0.                                  \tag{F47.1}
\]
Indeed, the finite presentation of $A$ and a finite set of generators
of $L/k$ allow a number field $F_0$ such that
$L\subset F_0\mathbb Q^{\mathrm{ab}}$ and $A$ is obtained by base
change from an abelian variety over $F_0$. Tsujimura's
Theorem~3.1 \cite{Tsujimura2020} says exactly that the Kummer map
on $A(F_0\mathbb Q^{\mathrm{ab}})$ is injective. Its kernel is the
intersection of the multiplication-by-$n$ images, so the injection
$A(L)\hookrightarrow A(F_0\mathbb Q^{\mathrm{ab}})$ proves
(F47.1). The compatible total Tate module of $A(L)$ is also zero:
each coordinate of a compatible torsion system is divisible by every
integer, and (F47.1) kills it. Separately, the height argument in the
rational-source lemma gives
\[
 D(L^*):=\bigcap_n(L^*)^n\subset\mu_\infty.                  \tag{F47.2}
\]
Neither assertion requires that $L$ have a discrete valuation.

\emph{A smooth-curve Kummer specialization lemma.}
Let $C/L$ be a smooth geometrically connected affine curve and let
\[
 \xi\in\varprojlim_n H^1_{\mathrm{et}}(C,\mu_n).
\]
Assume that for every finite $L'/L$ and every $z\in C(L')$ the
specialization $z^*\xi$ is the Kummer class of an actual element of
$(L')^*$. After a finite extension of $L$ there is an actual regular
unit $g$ on $C$ such that $\xi=[g]$. This extension is only an
extension of constants used in the proof; no descent of the given
topos morphism is assumed.

To prove the lemma, let $\overline C$ be the smooth proper completion
and $S=\overline C\setminus C$. Enlarge $L$ so that the points of
$S$ and one point $z_0\in C$ are rational. Purity for a smooth curve
and the Kummer exact sequence give compatible residue maps
\[
 H^1(C,\mu_n)\longrightarrow
       \bigoplus_{s\in S}\mathbb Z/n,
 \qquad
 \operatorname{res}(\xi)=(\lambda_s)_{s\in S},
 \quad\lambda_s\in\widehat{\mathbb Z}.                    \tag{F47.3}
\]
We first prove that every $\lambda_s$ is an integer.

Fix $s\in S$. By Riemann--Roch and interpolation of values and a
first derivative at finitely many points, there is a nonconstant
rational function
\[
 \pi:\overline C\longrightarrow\mathbb P^1_L
\]
which has a simple zero at $s$ and takes nonzero values at every
other point of $S$. For example allow a sufficiently large pole
divisor at a point different from $S$, prescribe the first-order
jet of a uniformizer at $s$, and prescribe the value $1$ at the
other boundary points. The map is finite. Remove from $\mathbb P^1$
the values $0,\infty$, its branch values, and $\pi(S)$, and enlarge
$L$ once more to split this finite set if necessary. This gives
\[
 W=\pi^{-1}(V)\subset C,\qquad W\longrightarrow V
 \quad\hbox{finite etale},
\]
where $V$ is a punctured affine line. Apply the compatible norm
maps to obtain
\[
 \eta=\operatorname{Cor}_{W/V}(\xi|_W)
        \in\varprojlim_n H^1(V,\mu_n).
\]
The specialization property is preserved. If a fibre above
$a\in V(L')$ is the disjoint union of $\operatorname{Spec}L_i$,
etale base change for the norm and the usual Kummer norm identity
give the Kummer class of
\[
        \prod_i N_{L_i/L'}(u_i)\in(L')^*,
\]
where $u_i$ realizes the evaluation of $\xi$ at the corresponding
point of $W$. Thus the rational-source integrality lemma applies to
$\eta$.

The residue at $0$ of a corestriction is
$\sum_{z\mid0}[\kappa(z):L]\operatorname{res}_z(\xi)$.
This is the divisor-of-norm formula, and can also be read directly
from the norm of a local parameter in the Kummer boundary map.
Points of $C$ contribute zero because $\xi$ extends over them;
among the points of $S$ only the rational point $s$ lies over $0$.
Consequently this residue is exactly $\lambda_s$, with coefficient
one. The lemma on $V$ makes it integral. The choice of a simple
zero also makes explicit that no ramification factor has been
silently inserted.

Put $D=\sum_s\lambda_s[s]$, an integral divisor. The next arrow in
the localization sequence sends the residue vector to the Kummer
Chern class of $D$. Therefore
\[
 c_1(D)=0\quad\hbox{in }H^2(\overline C,\mu_n)
       \quad\hbox{for every }n,
\]
and the Kummer injection
$\operatorname{Pic}(\overline C)/n\hookrightarrow
H^2(\overline C,\mu_n)$ implies
$[D]\in\bigcap_n n\operatorname{Pic}(\overline C)$.
Its degree is zero. Any divisor class $E$ with $nE=[D]$ likewise
has degree zero. The rational point $z_0$ identifies degree-zero
line bundles with $J(L)$ for the Jacobian $J$ and removes the
Brauer obstruction to representing a Picard-scheme point by a line
bundle. Equation (F47.1) therefore gives $[D]=0$.
Choose $g_0\in L(C)^*$ with $\operatorname{div}(g_0)=D$;
it is a unit on $C$.

All residues of $\xi-[g_0]$ now vanish. For every $n$ the map
$H^1(\overline C,\mu_n)\to H^1(C,\mu_n)$ is injective: purity
identifies its preceding support group with a negative-degree
cohomology group of $S$. Thus its residue-free lifts are unique,
and hence compatible as $n$ varies. The Kummer sequences give
\[
 0\longrightarrow\widehat{L^*}\longrightarrow
  \varprojlim_n H^1(\overline C,\mu_n)
  \longrightarrow T\operatorname{Pic}(\overline C)
  \longrightarrow0.
\]
Here the left inverse system has surjective transition maps, so
there is no $\varprojlim^1$ obstruction. The total Tate module on
the right is zero by (F47.1). Thus $\xi-[g_0]$ is a constant
completed Kummer class. Evaluate at $z_0$: the hypothesis identifies
it with the Kummer class of an actual constant $c\in L^*$.
Then $g=cg_0$ proves the lemma.

\emph{Algebraicity of the point function for a smooth curve.}
Let
$\phi:C_{\mathrm{et}}\to(\mathbb G_m)_{\mathrm{et}}$
be pinned and relative to $B G_L$, and write
$f:C(K)\to K^*$ for its point function. It is Galois equivariant
and sends $L'$-points to $L'$-points for every finite extension,
as checked in F.46. Its pulled-back Kummer class satisfies the
specialization lemma. After a finite extension choose
$g_0\in\mathcal O(C)^*$ representing it. Equation (F47.2) gives
\[
 f(x)/g_0(x)\in\mu_\infty\qquad(x\in C(K)).                 \tag{F47.4}
\]
If the generic image is a closed point, the inverse-image of its
open complement is empty by continuity and the Jacobson property.
The morphism factors through that closed subtopos. After splitting
its residue field it has a constant point function. We may therefore
assume the generic image is generic.

For each $b\in K$, base change to contain $b$ and pull back the
open $\mathbb G_m\setminus\{b\}$. It is a dense open
$C_b\subset C$ whose omitted geometric points are finite, and
$x\in C_b(K)$ exactly when $f(x)\ne b$. Applying the curve lemma
to the Kummer classes of $t-b$ gives a rational function $g_b$, a
unit on $C_b$, with
\[
 f(x)-b=d_b(x)g_b(x),\qquad d_b(x)\in\mu_\infty,
       \quad x\in C_b(K).                                 \tag{F47.5}
\]
Use the fixed proper field $M=F\mathbb Q^{\mathrm{ab}}\subset K$.
In any one argument, a finite extension $N/M$ contains the finitely
many coefficients needed there; the field $M$ itself stays fixed.
All the coefficients
$d_b(x)$ lie in $M^*$. The extension $K/M$ has infinite degree,
as does $K/N$ for every finite extension $N/M$ in $K$.
For clarity, the latter assertion can also be seen from the height
bound: if $K=N$ were finite over $M$, that bound would be uniformly
positive on its nontorsion elements, whereas $K$ contains all roots
of the rational number $2$.

The finite-exception quasi-rationality argument in F.46 applies to
an arbitrary smooth curve, with the following explicit replacement
for its polynomial-degree step. Descend the curve and the finitely
many rational functions under consideration to a finite extension
$N/M$. The poles of any fixed linear combination of two such
functions lie in one fixed divisor. If the combination is nonzero,
each zero has residue degree over $N$ bounded by the degree of that
divisor. Points of arbitrarily large residue degree exist: choose a
finite map of the proper curve to $\mathbb P^1_N$, choose values of
arbitrarily large degree over $N$, and take points above them.
They avoid any prescribed finite exceptional set. Thus a vanishing
at such a point forces the same rational-function identity as in
F.46. The two shifted outside-$M$ sets cover the cofinite domain,
and have infinite intersection by this same degree bound.

If $g_0$ is nonconstant, this proves that $f$ equals a rational
function $R\in K(C)$ on a cofinite subset. If $g_0$ is constant,
divide $f$ by one of its values and use (F47.4) to put its values
in $\mu_\infty\subset M$. Choose $b\notin M$. If $g_b$ were
nonconstant, (F47.5) would put its cofinite image in the vector
space $M+Mb$, which cannot contain a cofinite subset of $K$.
Hence $g_b$ is constant. Two different values of $f$ would, by
subtracting their equations, force this constant into $M$, and
then force $b\in M$. Thus $f$ is constant in this case.

We remove the finite exceptional set in the nonconstant case.
For every $b$, the rational function $(R-b)/g_b$ takes its values
in $\mu_\infty$ on a cofinite set. A nonconstant rational function
on a smooth curve has cofinite image in $K$, so this quotient is a
constant $c_b\in\mu_\infty$. Fix $x\in C(K)$ and choose
$b\ne f(x)$. Since $g_b$ is a unit at $x$, the identity
$R-b=c_b g_b$ first excludes a pole of $R$ at $x$. For every
$b\notin\{f(x),R(x)\}$, it then gives
\[
 \frac{f(x)-b}{R(x)-b}\in\mu_\infty.
\]
If $f(x)\ne R(x)$ this fractional linear expression takes all
values of $K\setminus\{0,1\}$ as $b$ varies, a contradiction.
Therefore $f=R$ on every geometric point. It follows that $R$ is
a regular unit on $C_K$. Galois equivariance and equality on the
Zariski-dense geometric points imply that $R$ descends to a regular
unit on the original curve over $L$. This proves the full
point-function assertion for every smooth affine curve and every
finite extension of the base field.

It also gives the equality of compatible Kummer classes that will
be used below. Over the finite extension $L'/L$ employed in the
specialization lemma, the rational function $R/g_0$ has values in
$\mu_\infty$ at every geometric point and is therefore a constant
$\zeta\in(L')^*$. Evaluation at a rational point of the curve gives
$\zeta\in D((L')^*)$, so $[R]=[g_0]=\xi$ over $L'$.
If $C(L)$ is nonempty, equality already holds over $L$. Indeed,
for every $n$ the residual class in $H^1(C,\mu_n)$ is geometrically
zero, and hence is a constant class by the low-degree
Hochschild--Serre sequence. Its evaluation at an $L$-point is zero,
because $R$ is the actual point function there. The same sequence,
or the rational-point section, shows that this constant class is
zero. Without an $L$-point, equality after a finite extension is all
that is asserted or needed for the geometric residue tests.
